% Generated from locale/tr/content/second-order-logic/metatheory/loewenheim-skolem.tex; do not hand-edit.


\olfileid[tr]{sol}{met}{lst}

\olsection{Löwenheim--Skolem Teoremi İkinci Dereceden Mantıkta Geçersizdir}

\begin{explain}
(Aşağı Doğru) Löwenheim--Skolem Teoremi, sonsuz bir modeli olan her
!!{sentence} kümesinin !!a{enumerable} modeli bulunduğunu söyler. Bu da tamlık
teoreminin bir sonucudur: tamlık ispatı, her tutarlı !!{sentence} kümesi için
bir model üretir ve bu model !!{enumerable}{}dir. Bir de bir !!{sentence}
kümesinin !!a{denumerable} modeli varsa !!a{nonenumerable} modelinin de
bulunduğunu güvence altına alan Yukarı Doğru Löwenheim--Skolem Teoremi
vardır. İki teorem de ikinci dereceden mantıkta geçersizdir.
\end{explain}


\begin{thm}
\ollabel{thm:sol-no-ls} Löwenheim--Skolem Teoremi ikinci dereceden mantıkta
geçersizdir: !!{sentence}s vardır ki sonsuz modelleri bulunur, fakat
!!{enumerable} modelleri bulunmaz.
\end{thm}

\begin{proof}
\[
\fn{Count} \ident \lexists[z][\lexists[u][\lforall[X][((X(z) \land
      \lforall[x][(X(x) \lif X(u(x)))]) \lif \lforall[x][X(x)])]]]
\]
ifadesinin !!a{structure} $\Struct{M}$ içinde ancak ve ancak $\Domain{M}$
!!{enumerable} ise doğru olduğunu anımsayın. Dolayısıyla
$\lnot\fn{Count}$, $\Struct{M}$ içinde ancak ve ancak $\Domain{M}$
!!{nonenumerable} ise doğrudur. Böyle !!{structure}s vardır: herhangi bir
!!{nonenumerable} kümeyi !!{domain} olarak, örneğin $\Pow{\Nat}$ ya da
$\Real$'i alın. Böylece $\lnot\fn{Count}$'un sonsuz modelleri vardır, fakat
!!{enumerable} modeli yoktur.
\end{proof}

\begin{thm}
!!{sentence}s vardır ki !!{denumerable} modelleri bulunur, fakat
!!{nonenumerable} modelleri bulunmaz.
\end{thm}

\begin{proof}
$\fn{Count}\land\fn{Inf}$, $\Nat$ içinde doğrudur; fakat
!!{structure}~$\Struct{M}$, $\Domain{M}$ !!{nonenumerable} olduğunda bu
tümceyi sağlamaz.
\end{proof}


